\documentclass{article}

% if you need to pass options to natbib, use, e.g.:
%     \PassOptionsToPackage{numbers, compress}{natbib}
% before loading neurips_2026

% The authors should use one of these tracks.
% Before accepting by the NeurIPS conference, select one of the options below.
% 0. "default" for submission
\usepackage{neurips_2026}
% the "default" option is equal to the "main" option, which is used for the Main Track with double-blind reviewing.
% 1. "main" option is used for the Main Track
%  \usepackage[main]{neurips_2026}
% 2. "position" option is used for the Position Paper Track
%  \usepackage[position]{neurips_2026}
% 3. "eandd" option is used for the Evaluations & Datasets Track
% \usepackage[eandd]{neurips_2026}
 % if you need to opt-in for a single-blind submission in the E&D track:
 %\usepackage[eandd, nonanonymous]{neurips_2026}
% 4. "creativeai" option is used for the Creative AI Track
%  \usepackage[creativeai]{neurips_2026}
% 5. "sglblindworkshop" option is used for the Workshop with single-blind reviewing
 % \usepackage[sglblindworkshop]{neurips_2026}
% 6. "dblblindworkshop" option is used for the Workshop with double-blind reviewing
%  \usepackage[dblblindworkshop]{neurips_2026}

% After being accepted, the authors should add "final" behind the track to compile a camera-ready version.
% 1. Main Track
 % \usepackage[main, final]{neurips_2026}
% 2. Position Paper Track
%  \usepackage[position, final]{neurips_2026}
% 3. Evaluations & Datasets Track
% \usepackage[eandd, final]{neurips_2026}
% 4. Creative AI Track
%  \usepackage[creativeai, final]{neurips_2026}
% 5. Workshop with single-blind reviewing
%  \usepackage[sglblindworkshop, final]{neurips_2026}
% 6. Workshop with double-blind reviewing
%  \usepackage[dblblindworkshop, final]{neurips_2026}
% Note. For the workshop paper template, both \title{} and \workshoptitle{} are required, with the former indicating the paper title shown in the title and the latter indicating the workshop title displayed in the footnote.
% For workshops (5., 6.), the authors should add the name of the workshop, "\workshoptitle" command is used to set the workshop title.
% \workshoptitle{WORKSHOP TITLE}

% "preprint" option is used for arXiv or other preprint submissions
 % \usepackage[preprint]{neurips_2026}

% to avoid loading the natbib package, add option nonatbib:
%    \usepackage[nonatbib]{neurips_2026}

\usepackage[utf8]{inputenc} % allow utf-8 input
\usepackage[T1]{fontenc}    % use 8-bit T1 fonts
\usepackage{hyperref}       % hyperlinks
\usepackage{url}            % simple URL typesetting
\usepackage{booktabs}       % professional-quality tables
\usepackage{amsfonts}       % blackboard math symbols
\usepackage{nicefrac}       % compact symbols for 1/2, etc.
\usepackage{microtype}      % microtypography
\usepackage{xcolor}         % colors

% Note. For the workshop paper template, both \title{} and \workshoptitle{} are required, with the former indicating the paper title shown in the title and the latter indicating the workshop title displayed in the footnote. 
\title{INSTANT - INdia Satellite Temporal Analysis for Nowcasting Trends.}


% The \author macro works with any number of authors. There are two commands
% used to separate the names and addresses of multiple authors: \And and \AND.
%
% Using \And between authors leaves it to LaTeX to determine where to break the
% lines. Using \AND forces a line break at that point. So, if LaTeX puts 3 of 4
% authors names on the first line, and the last on the second line, try using
% \AND instead of \And before the third author name.


\author{%
  Shreeyans Arora \\
  Department of Computer Science\\
  Indian Insitute of Information Technology, Nagpur\\
  \texttt{bt32csa005@iiitn.ac.in} \\
  % examples of more authors
  \And
  Muskaan Maurya \\
  Independent Researcher \\
  \texttt{muskaan.maurya06@gmail.com} \\
  \And
  Utkarsh Mall \\
  Mohamed bin Zayed University of Artificial Intelligence \\
  \texttt{utkarsh.Mall@mbzuai.ac.ae} \\
  % \And
  % Coauthor \\
  % Affiliation \\
  % Address \\
  % \texttt{email} \\
  % \And
  % Coauthor \\
  % Affiliation \\
  % Address \\
  % \texttt{email} \\
}


\begin{document}


\maketitle


\begin{abstract}
  Precipitation nowcasting has gained significant importance over past few years due its relevance in short-term weather forecasting. 
  Several datasets have been released inorder to effectively carry out training and testing of deep-learning based nowcast models. 
  However, past works have a gap in nowcast research on Indian subcontinent speacially due unavailibility of high resolution, curated datasets. 
  This paper introduces a new benchmark dataset named INSTANT build to carry out studies on precipitation on 4 major Indian regions namely North (Jammu and Kashmir), North-Eastearn (7 sisters), Eastern (Jharkhand, Orissa and West Bengal) and South Indian (Kerala, Tamil Nadu). 
  INSTANT is a satellite based, spatially and temporally aligned dataset created from raw weather data of MOSDAC which was collected using the IMAGER sensor and INSAT-3DR satellite. 
  It consists of 320472 precipitaion images each of 896km x 896km  and timestamp of half hour.
  We have performed a comparative analysis of three deep learning based nowcast models on our proposed dataset which serves as a baseline. Moreover, we have also introduced a new metric to measure the performance of models on basis of their behaviour on long-lead time. 
  The paper further discusses the collection methodology along with other evaluation metrics.
\end{abstract}



\section{Introduction}
\begin{itemize}
  \item Accurate short-term forecasting is essential for predicting the storm movements to prevent the damages caused by it. People lose their livelihood due to these calamities which, in most cases, have irreversible losses. 
Precipitation nowcasting can provide accurate short-term precipitation results due its ability to provide highly detailed predictions of localized weather phenomena that undergo significant changes.
  \item This had lead to an increase in demand for the use of precipitation nowcast for forecasting which made researchers exapnd their scope from traditional statistical method to DL-based techniques.cite. (write about PINNs in model discusssion)\\
Training and evaluating and these models require leveraging the use of rich data sources like GOES and NEXRAD but due their large size, direct use becomes difficult. 
  \item  SEVIR, cite, has tried to deal with this problem by significantly downsampling these data and also cmobining them for different modatlities increase their usage domain. But region is problem, GOES-16 and GOES-17 and NAXRAD are only for US region.
  \item  There are few publicly available weather dataset for India. 
IMDAA Indian Monsoon Data Assimilation and Analysis: Regional atmospheric reanalysis; used for studying monsoon variability, extreme weather diagnostics, and climate trends.
IndiaWeatherBench: Machine learning-ready regional weather forecasting; used for benchmarking deep learning models on Indian weather variables.
IPED Indian Precipitation Ensemble Dataset: High-resolution daily gridded precipitation; used for hydrological modeling, rainfall trend analysis, and model verification.
IMD Gridded Temperature/Rainfall: Historical climate analysis; used for long-term temperature trends and seasonal rainfall monitoring across a 1.0° x 1.0° or 0.25° x 0.25° grid.
Although good for tasks like Predicting regional temperature, wind speed, and precipitation amounts over the next several days, none of them are good for short-term forecasting task. 
  \item To help address this issue, we introduce INdia Satellite Temporal Analysis for Nowcasting Trend (INSTANT), designed for advancing research on Indian nowcasting. Brief about INSTANT (13500 rainfall events)


\end{itemize}

\section{INSTANT - India Satellite Temporal Analysis for Nowcasting Trends}
\begin{itemize}
  \item INSTANT is a satellite-based, spatially and temporally aligned dataset created from raw weather data of MOSDAC on Indian subcontinent which was collected using the IMAGER sensor and INSAT-3DR satellite. MOSDAC used INSAT multispectral Rainfall Algorithm Technique (IMSRA) to retrieve precipitation information from raw staellite images. This algorithm estimates precipitation based on non-linear power law relationship established between infrared (IR) brightness temperatures (Tbs, 10.7 µm observations) and TRMM-PR surface rain rate.
  \item (original image info) image type is infrared. Origianl image is of shape []  where each pixel is ~4km and each channel represents. coverage - Max Long.	120.0, Min Long.	30.0, Max Lat.	40.0,Min Lat.	40.0. Available from date 08/11/2016 to present.
  \item (our changes) extracted data for 4 years (01/01/2021 - 31/12/2024) divided into four region with following coordinates:\\
  Northern:  Max Long.	81.0, Min Long.	73.0, Max Lat. 37.0, Min Lat. 29.0\\
  Northeastern:  Max Long.	98.0, Min Long.	90.0, Max Lat. 30.0, Min Lat.	22.0\\
  Easthern:  Max Long.	90.0, Min Long.	82.0, Max Lat.	26.0,Min Lat.	18.0\\
  southern india:  Max Long.	81.0, Min Long.	73.0, Max Lat.	15.0,Min Lat.	7.0\\
  each image per region is of shape 112 x 112. the original image (4km res) showed gridded artifact in precipitation mapping (a problem from MOSDAC only). The gridded pattern at 4km is a retrieval artifact where the algorithm assigns uniform rainfall values to fixed spatial blocks, making the data look artificially blocky rather than smooth.It is important to treat this artifact as, the model may learn to reproduce the artificial block pattern instead of actual rainfall dynamics, leading to predictions that mimic the sensor artifact rather than true precipitation structure. At 8km, each pixel naturally covers a larger area so the block boundaries blend together, producing a smoother and more physically realistic rainfall field. Now each image is of size 896km x 896km.
  \item total images: 320472, each of HDF5 format and shape [3, 112, 112] total disk space of .
\end{itemize}

\subsection{Event selection}
location graph, month graph (histogram - dataset count vs month)

\section{Training and Evaluations} 

\subsection{Model Architecture and Loss function}
\begin{itemize}
  \item \textbf{Unet}(cite): The proposed model follows a U-Net-based encoder–decoder design with a 12-channel input and an 18-channel output at a spatial resolution of 112×112. The encoder begins with an initial double convolution block followed by four successive downsampling stages, each comprising a 2×2 max-pooling operation and a double convolution block consisting of two sequential 3×3 convolutions, each followed by batch normalization and ReLU activation, 
  progressively expanding the feature channels from 64 to 512. The bottleneck operates at 512 channels, after which the decoder symmetrically reconstructs the spatial resolution through four upsampling stages using bilinear interpolation, halving the channel count at each stage while incorporating skip connections from the corresponding encoder levels via concatenation to preserve spatial detail; a padding mechanism ensures spatial compatibility between encoder and decoder 
  feature maps during concatenation. A final 1×1 convolution projects the resulting 64-channel feature map to the 18-channel output.
  \item \textbf{Earthformer}(cite): The model follows a Cuboid Transformer encoder–decoder design operating on channels-last spatiotemporal tensors of shape (B, T, H, W, C), taking 4 input frames of spatial resolution 112×112 with 3 channels and predicting 6 future frames of the same resolution. Prior to the transformer stages, the input is passed through a stacked convolutional patch-merging encoder comprising 3 successive layers with progressive downscaling factors
   of (3, 2, 2) and channel dimensions expanding through [32, 64, 128], each layer applying 2 convolutional operations with Leaky ReLU activation; a symmetric stacked upsampling decoder with transposed convolutions reverses this process in the decoder. Additive spatiotemporal positional embeddings are injected after the initial downsampling via separate temporal, height, and width learned embeddings (t+h+w scheme). The encoder and decoder each consist of 2 Cuboid Transformer blocks employing axial self-attention, 
   which decomposes full spatiotemporal attention into three independent 1D attention passes along the temporal, height, and width axes respectively, with 4 attention heads, relative positional biases, dropout rates of 0.1 for attention, projection, and feed-forward layers, and inter-block feed-forward networks with GELU activation and layer normalization. A set of 8 learnable global vectors is incorporated in the encoder to capture long-range global context via dedicated global self-attention. The decoder additionally applies 
   1×1 cuboid cross-attention to condition each decoder block on the corresponding encoder memory, with the decoder query sequence initialized from zeros and projected via a learned linear layer. A final linear projection maps the decoder output back to the 3-channel target space.
  \item \textbf{NowcastNet}(cite):The NowcastNet model follows a Generative Adversarial Network (GAN) framework comprising a generator and a temporal discriminator. The generator accepts a concatenated input of shape $(B, C_{in}, H, W)$, where $C_{in} = 10$ corresponds 
to 4 context frames and 6 evolution prediction frames at $112 \times 112$ spatial 
resolution. The generative encoder applies three successive $\texttt{MaxPool2d}(2)$ 
operations to downsample the input to a $14 \times 14$ feature map of dimension 
$n_{gf} \times 8$ channels. A noise branch projects a latent noise vector $z$ of 
dimension 32 through a linear projector followed by a depth-to-space (PixelShuffle) 
rearrangement, then spatially resized via adaptive average pooling to match the 
encoder output; the noise feature and encoded feature are concatenated and passed 
to the generative decoder, which reconstructs 6 predicted frames of shape 
$(B, 6, H, W)$. The temporal discriminator processes the concatenated real or 
generated sequence through three parallel branches --- a 2D spectral-norm 
convolution ($9 \times 9$, stride 2) and two 3D spectral-norm convolutions 
operating over the temporal dimension with short-range (kernel $(4, 9, 9)$, 
stride $(1, 2, 2)$) and long-range (kernel $(4, 9, 9)$, stride $(4, 2, 2)$) 
receptive fields respectively --- whose outputs are concatenated and passed 
through four successive L2 blocks with spectral-norm convolutions, batch 
normalization, and Leaky ReLU activations, followed by a spectral-norm linear 
layer producing raw logits of shape $(B, 1)$.

\begin{equation}
    \mathcal{L}_{D} = -\mathbb{E}\left[\log \sigma(D(x, Y))\right] - \mathbb{E}\left[\log\left(1 - \sigma(D(x, \hat{Y}))\right)\right]
\end{equation}

\begin{equation}
    \mathcal{L}_{G} = -\mathbb{E}\left[\log \sigma(D(x, \hat{Y}))\right]
\end{equation}

where $\sigma(\cdot)$ denotes the sigmoid function, $D(\cdot)$ denotes the temporal discriminator 
outputting raw logits, $G(\cdot)$ denotes the generator, $x \in \mathbb{R}^{B \times C_{in} \times H \times W}$ 
is the concatenated context and evolution prediction input, and $z \in \mathbb{R}^{B \times d_z \times h \times w}$ 
is a sampled noise vector with latent dimension $d_z = 32$.
\end{itemize}

Unet amd Earthformer models are trained using the Mean Absolute Error (MAE) loss, also known as the $\ell_1$ loss. Given a predicted output $\hat{Y} \in \mathbb{R}^{B \times T_{out} \times H \times W \times C}$ and the corresponding ground truth $Y \in \mathbb{R}^{B \times T_{out} \times H \times W \times C}$, the loss is defined as:

\begin{equation}
    \mathcal{L} = \frac{1}{B \cdot T_{out} \cdot H \cdot W \cdot C} \sum_{b=1}^{B} \sum_{t=1}^{T_{out}} \sum_{h=1}^{H} \sum_{w=1}^{W} \sum_{c=1}^{C} \left| Y_{b,t,h,w,c} - \hat{Y}_{b,t,h,w,c} \right|
\end{equation}
where $B$ is the batch size, $T_{out}$ is the number of predicted output frames ($T_{out} = 6$), $H$ and $W$ are the spatial dimensions ($H = W = 112$), and $C$ is the number of channels ($C = 3$). Both the U-Net and Cuboid Transformer models minimize $\mathcal{L}$ using the Adam optimizer.



\subsection{Long-lead time analysis}
input output discuss

\subsection{Evaluation Metric}
All metric discussion with formulas

\subsection{Satellite nowcast Prediction}
model prediction images
results table

\section{Conclusion}

\section{Broader impact}

\section{Acknowledgement and funding}


\section*{References}


References follow the acknowledgments in the camera-ready paper. Use unnumbered first-level heading for
the references. Any choice of citation style is acceptable as long as you are
consistent. It is permissible to reduce the font size to \verb+small+ (9 point)
when listing the references.
Note that the Reference section does not count towards the page limit.
\medskip


{
\small


[1] Alexander, J.A.\ \& Mozer, M.C.\ (1995) Template-based algorithms for
connectionist rule extraction. In G.\ Tesauro, D.S.\ Touretzky and T.K.\ Leen
(eds.), {\it Advances in Neural Information Processing Systems 7},
pp.\ 609--616. Cambridge, MA: MIT Press.


[2] Bower, J.M.\ \& Beeman, D.\ (1995) {\it The Book of GENESIS: Exploring
  Realistic Neural Models with the GEneral NEural SImulation System.}  New York:
TELOS/Springer--Verlag.


[3] Hasselmo, M.E., Schnell, E.\ \& Barkai, E.\ (1995) Dynamics of learning and
recall at excitatory recurrent synapses and cholinergic modulation in rat
hippocampal region CA3. {\it Journal of Neuroscience} {\bf 15}(7):5249-5262.
}


\newpage
\input{checklist.tex}


\end{document}